\song{Hej teto (Michal Horák)}
\begin{multicols}{2}

\re{}
\textit{capo 3}

\ve{}
Měl jsem \ch{Ami}hlídat tetě \ch{C}křečka,\\                  
\ch{G}dala mi ho se slovy,\\
ať \ch{D}u mě chvíli přečká.
   
\ve{}
Ale křečka uhlídati\\
nelze jen tak lehce,\\
obzvlášť když se vám moc nechce.

\ve{}
Po pokoji ať si běhá,\\
na gauč jsem si lehnul,\\
vzápětí se ani nehnul.

\ve{}
Tlačilo mě do zad kromě\\
polštářů i cosi,\\
už se vidim tetu prosit

\cho{}
Hej \ch{F}teto! Já ho \ch{C}zabil.\\
\ch{E7}Už se to nestane,\\
jo, \ch{Ami}ten už asi nevstane.\\
Hej \ch{F}teto! Nekřič \ch{C}na mě.\\
\ch{E7}Já ho prostě neviděl,\\
\ch{Ami}už dost jsem se nastyděl.\\
Hej \ch{F}teto! \ch{C \ E7 \ Ami \ F \ C \ E7}

\columnbreak


\ve{}
Tak jsem kouknul, co mě tlačí,\\
seděl jsem jenom\\
na televizním ovladači.

\ve{}
Ale křeččí švitoření\\
slyšet není, jenom zvláštní smrad\\
se line od topení.

\ve{}
Zase se mi dech zatajil,\\
spečeného křečka bych\\
před tetou neobhájil.

\ve{}
A dřív než tam sebou seknu,\\
měl bych asi vymyslet,\\
jak tetě potom řeknu.

\cho{}

\cho{}
Hej teto! Já ho fakt zabil.\\
Ten bídák sežral magnet\\
a k topení se přitavil.\\
Hej teto, co na to říci?\\
Můžeš si ho maximálně\\
připnout na lednici. \cho{}

\end{multicols}
\esong{}